\section{Introdução}
\label{sec:introducao}
Monitorar a atividade da Ethereum envolve selecionar, entre milhões de operações diárias, transações que merecem investigação após sua execução. Transferências, contratos inteligentes e aplicações financeiras compartilham os mesmos blocos, mas diferem em consumo de gás, taxas, posição e recorrência dos participantes. Esses atributos permitem caracterizar desvios e organizar a triagem em um livro-razão distribuído. A ordenação também possui valor econômico: Daian et al.\ \cite{Daian2020} documentam bots de arbitragem e competição por prioridade, e Torres et al.\ \cite{Torres2021} investigam padrões de \textit{front-running}. Escores de anomalia podem apoiar a escolha de operações para análise, articulando caracterização da rede e esforço de investigação.

O problema considerado é compreender como dois mecanismos não supervisionados de detecção respondem às mudanças temporais dessa atividade. Aplicar um corte fixo não garante um volume estável de sinalizações; combinar detectores também não informa, por si só, quais transações cada método acrescenta. Além disso, uma boa ordenação global pode coexistir com baixa concentração de operações de interesse nas primeiras posições do ranking. Essas diferenças são relevantes quando a capacidade de inspeção é restrita e a distribuição transacional muda ao longo do tempo.

Comparamos Autoencoder (AE) e Isolation Forest (IF) em 523.937.952 transações disponíveis na fonte AWS nos 365 dias de 2025, descritas pelos mesmos 21 atributos \textit{on-chain}. Os modelos são ajustados de janeiro a abril, com hiperparâmetros escolhidos em estudo anterior, e os limiares são calibrados em junho e mantidos fixos. Investigamos três perguntas: \textbf{RQ1}, como o volume sinalizado responde aos diferentes recortes sob cortes congelados? \textbf{RQ2}, quanto os modelos concordam e quais são suas contribuições exclusivas? \textbf{RQ3}, como a ordenação global se relaciona à concentração de transações associadas a bots MEV sob um mesmo orçamento de análise?

A contribuição é uma caracterização temporal e operacional desses mecanismos na mesma população transacional. Farrugia et al.\ \cite{Farrugia2020} classificam contas com atributos agregados, enquanto a literatura de \textit{front-running} identifica relações estruturais entre operações. Aqui, a unidade é a transação executada; a comparação relaciona mudanças marginais, conjuntos compartilhados e exclusivos e decisões sob orçamento fixo. A associação histórica do remetente no Etherscan Label Cloud fornece uma referência auxiliar após a inferência. O estudo quantifica como as prioridades produzidas por reconstrução e isolamento divergem ao longo de 2025, sem tratar essa associação como confirmação de ataques individuais.

A Seção~\ref{sec:relacionados} situa essa comparação na literatura. A Seção~\ref{sec:metodologia} apresenta os dados e o protocolo; a Seção~\ref{sec:resultados} responde às perguntas por meio das taxas, interseções e rankings. A Seção~\ref{sec:conclusoes} sintetiza os achados e a continuidade da investigação.
